\documentclass{report}
\usepackage{amsmath,amssymb}
\setlength{\parindent}{0mm}
\setlength{\parskip}{1em}
\begin{document}
\begin{center}
\rule{6in}{1pt} \
{\large Edmond Chow \\
{\bf Reducing communication costs in distributed relaxation methods}}

School of Computational Science and Engineering \\ Georgia Institute of Technology \\ Atlanta \\ GA \\ USA
\\
{\tt echow@cc.gatech.edu}\\
Jordi Wolfson-Pou\end{center}

Inter-node communication has high energy cost compared to computation.
In exascale computing, reducing communication and thus energy cost may
become even more important than reducing overall time. We apply this
extreme objective to relaxation-type iterative methods and propose a new,
``parallel Southwell'' iteration for solving sparse linear systems.
Standard Southwell iteration relaxes one row at a time, corresponding
to the row with the highest residual, and is a sequential algorithm by
definition. Southwell generally converges faster than Gauss-Seidel in
terms of the number of relaxations. Since each relaxation is associated with
communication, communication cost is also reduced. However, to achieve
this faster convergence, Southwell needs global communication after
every relaxation step to determine the row with the highest residual.

In parallel Southwell, we approximate the standard Southwell iteration
by 1) avoiding global communication, and 2) allowing some rows to be
relaxed simultaneously. To describe the method, assume one processor
node assigned to one row or grid point. At each iteration, each node
determines if it is the node with the largest residual among its
neighbors in the grid. If so, then this node relaxes its row and sends
updated values to its neighbors. Nodes that have recently relaxed
do not relax again for a given number of iterations, to prevent too
many nodes relaxing simultaneously. Open questions include how each
node maintains an estimate of the residual on neighboring nodes (this
process does not need to be exact). Also, the algorithm is ideally
implemented as an asynchronous method, and thus we discuss the challenges
of implementing this method using MPI one-sided communication (remote
memory access). The goal is to develop a method to solve problems with
less communication than with Gauss-Seidel or Jacobi with the possibility
of also being faster on extreme scale distributed systems.


\end{document}
